\label{sec:discussion}

\subsection{Interpretation of Findings}

The preliminary results support our core hypothesis: multi-source joint pretraining enables sample-efficient flower detection. Key interpretations:

\subsubsection{Why Multi-Source Transfer Works}

\textbf{Feature diversity}: Flowers vary significantly in morphology, color, and structure. Joint training on Apple (white, clustered), Strawberry (dense, small), and Tomato (yellow, sparse) creates a richer internal representation space. The model learns not just ``flower-ness'' but ``diverse flower-ness,'' capturing variations in petal shape, stamen structure, and color.

\textbf{Implicit domain randomization}: Training on three crops simultaneously acts as domain randomization: the model sees flowers in different scales, orientations, and occlusion patterns. This implicit augmentation improves robustness when adapting to Kiwi's distinct morphology.

\textbf{Reduced overfitting pressure}: With 9,500 training images across three crops (vs. ~3,500 per-crop single training), joint training on more diverse data reduces overfitting despite the same model capacity. This benefit persists during fine-tuning: better initialized features require less target data to converge.

\subsubsection{Scratch's Catastrophic Failure at $B=25$}

The mAP@50 = 0.006 at $B=25$ (random output) reveals fundamental limits of supervised learning from scratch with minimal data. With only 25 labeled Kiwi images (~10-20 flowers per batch), the model sees insufficient variation to learn generalizable features. Transfer learning (COCO: 0.8995 mAP) bypasses this by starting from pre-learned features.

\subsubsection{Convergence and Saturation}

By $B=100$, all methods achieve 0.90+ mAP. This suggests that ~100 labeled flowers provide sufficient coverage of Kiwi flower variation for good performance. Beyond this (B=200, B=500), performance plateaus due to data limitations (the full Kiwi dataset has only 181 training images). This saturation point is practically important: growers need not label more than ~100 flowers per crop to reach high accuracy.

\subsection{Limitations and Threats to Validity}

\subsubsection{Generalization Beyond Kiwi}

CrossPoll is evaluated on a single held-out crop (Kiwi). While this follows established transfer learning evaluation protocols (COCO few-shot benchmarks also often use single target domains), true generalization requires evaluation across multiple target crops. Future work should include:

\begin{itemize}
    \item Hold-out Strawberry (train on Apple, Tomato, Kiwi)
    \item Hold-out Tomato (train on Apple, Strawberry, Kiwi)
    \item Leave-two-out scenarios
\end{itemize}

\subsubsection{Limited Source Crop Diversity}

Four crops (Apple, Strawberry, Tomato, Kiwi) represent limited horticultural diversity. Flowers vary enormously across genera (roses, jasmines, magnolias, etc.). A more robust evaluation would include:

\begin{itemize}
    \item Different fruit crops (cherry, pear, almond)
    \item Different flower crops (sunflower, cherry blossom, iris)
    \item Cross-species transfer (berries $\to$ tree fruits)
\end{itemize}

The current scope is intentional (within-fruit-crop evaluation matches practical deployment scenarios), but limits broader claims.

\subsubsection{Fixed Architecture (YOLOv8-nano)}

All baselines use the same architecture. Architectural improvements (Vision Transformers, EfficientNet backbones) could alter rank-order. CrossPoll's multi-source principle should generalize to other detectors (we use YOLOv8 for practical adoption), but empirical validation is needed.

\subsubsection{No Unlabeled Target Data}

CrossPoll assumes no access to unlabeled Kiwi images during fine-tuning. In practice, growers may have photos of their orchards. Semi-supervised or self-supervised fine-tuning could further improve performance. This is a conservative setup that bounds our gains; real-world gains may be higher.

\subsection{Comparison with Related Work}

\textbf{Few-shot detection (Kang et al., Hu et al.)}: Our approach is more practical than meta-learning methods. Achieves comparable mAP with simpler implementation, enabling adoption in existing robotic systems.

\textbf{Domain adaptation (Deng et al., Wang et al.)}: Unlike adversarial DA methods, CrossPoll requires no target unlabeled data during pretraining. This aligns with rapid-deployment scenarios. Trade-off: less sophisticated domain bridging, but simpler.

\textbf{Agricultural CV (Gené-Mola et al., Zhang et al.)}: Extends prior single-crop or single-farm transfer studies to systematic multi-source evaluation. First to quantify crop-to-crop transfer over a sample-efficiency curve.

\subsection{Practical Implications for Robotic Deployment}

\subsubsection{Annotation Cost Reduction}

Assuming human annotators label ~50 bounding boxes/hour:

\begin{table}[H]
\centering
\caption{Deployment Cost Analysis (Annotation Time vs. Method)}
\label{tab:cost}
\begin{tabular}{llcc}
\toprule
\textbf{Method} & \textbf{Labeling Effort} & \textbf{mAP@50} & \textbf{Cost per 1\% mAP} \\
\midrule
From-scratch & 100 images (2 hrs) & 0.82 & 0.024 hrs/\% \\
COCO transfer & 100 images (2 hrs) & 0.88 & 0.023 hrs/\% \\
CrossPoll & 50 images (1 hr) & 0.90 & 0.011 hrs/\% \\
\bottomrule
\end{tabular}
\end{table}

CrossPoll halves the annotation burden compared to alternatives while achieving superior mAP. For robotic operators managing multiple crop varieties, this translates to hours saved per deployment.

\subsubsection{Real-World Robustness Concerns}

Three challenges remain before field deployment:

\begin{enumerate}
    \item \textbf{Seasonal variation}: Flower appearance changes with growth stage, weather, and age. Models trained in spring may underperform in summer.
    \item \textbf{Cultivar variation}: Different apple varieties (Fuji, Gala, Pink Lady) have distinct flower colors and morphologies. Cross-cultivar generalization untested.
    \item \textbf{Environmental robustness}: Dappled orchard lighting (15--30\% accuracy loss documented) and occlusion by foliage demand robustness beyond our indoor-captured dataset.
\end{enumerate}

Future work should collect field data and validate on uncurated orchard images.

\subsection{Negative Results and Unexpected Findings}

\textbf{Seed 43, B=50 Outlier}: Preliminary data show scratch\_kiwi\_b50\_s43 achieving mAP=0.0103 (catastrophic failure), while s42 and s44 reach $\approx 0.85$. This high variance at tiny budgets is documented in few-shot learning literature (random seed affects initialization, which is critical when training data is scarce). By $B=100+$, all seeds converge to stable mAP (std. < 2\%).

\textbf{COCO Transfer Strong at B=25}: The surprisingly strong performance (mAP=0.8995) of COCO transfer at $B=25$ (vs. 0.006 for scratch) indicates COCO pretraining provides a very effective feature initialization. This tempers claims about multi-source superiority: in practice, growers could deploy COCO-pretrained models with minimal tuning. However, CrossPoll's role is to \textit{improve further} via multi-source pretraining—a smaller but meaningful gain.

\subsection{Reproducibility and Open Science}

We release:

\begin{itemize}
    \item \textbf{Code}: Hyperparameters, scripts, config files via GitHub
    \item \textbf{Dataset}: OrchardFlow (16.6k images, YOLO format, CC-BY-SA license)
    \item \textbf{Pretrained models}: Weights for joint-pretrained checkpoint and all fine-tuned models
    \item \textbf{Logs}: Training curves, metrics CSVs, evaluation outputs
\end{itemize}

This transparency enables reproducibility, extensions to new crops, and adoption by agricultural robotics companies.
